% AI-written, AI cross-reviewed; no human sign-off; not a claim that RH is proved
% Manuscript v3, heading p. 17; original out/r23_host.tex, SHA-256 a416cbdc1776136bfb4e4aff18286f0a1e80ddfd07386b3043525371b2073f69.
\subsection{P015. Selecting a host and joining the estimates on one family}
\label{r23:host}

\begin{proposition}[The host for a primitive character]
Fix a primitive character $\psi$ of conductor $M$ and put $\Lambda=\operatorname{cyc}(\psi)$.
Assume $V\mathrm{WinAt}(\psi,401)$ as in \hyperref[sub:B07]{B07}.
Given $0<\varepsilon<1/52$, there are constants $C_{T2},C_8,c_6,C_b,x_0$
with $C_8\ge1$ and $c_6,C_b>0$, and a system of windows of exponent
$3/2$. At every member scale $s$ with $s.\varepsilon=\varepsilon$ and
$x=s.x\ge x_0$ there are ranks $R$, a host $H$ and a prime $q_0$ with
$\mathrm{RankVal}(s,R)$, $H.C_{T2}=C_{T2}$ and
$\mathrm{HostGood}(s,C_8,4,c_6,C_b,H,q_0)$. On this very host,
\[
  \mathrm{EAtH}(\Lambda,H,TT,x,\varepsilon,4)
  \quad\text{and}\quad
  \mathrm{VVle}(H.E,R.t,\sqrt{\exp(x^\varepsilon)})
\]
hold. The first predicate controls the kept characters, their coset shares
and the other frequencies. The second bounds all finite prime sums for
the indicated rank and every allowed choice of local data.
\end{proposition}
\noindent\nolinkurl{UpperTheory.Top.supplyChi}, \path{CharacterBoxes/SupplyChi.lean:8--16};
\nolinkurl{GChi}, \path{CharacterBoxes/Supply.lean:38--43};
\nolinkurl{HostGood}, \nolinkurl{HostExW},
\path{GlobalTwisting/ThinScalesIffHostBound.lean:101--123};
\nolinkurl{EAtH}, \path{ExplicitFormulaErrorOfHostConditions.lean:144--151};
\nolinkurl{VVle}, \path{CharacterBoxes/ExistsPageN.lean:12--18}.

\begin{formalstatement}[Inputs to the host selection]
Inside \decid{D510}, the prime-pair estimate gives two successive
positive proportions of good primes, at sufficiently large $Y$:
put $z=Y^{\varepsilon/5}$, $z'=z(\log Y)^2$, $X=Y^J$ with $J\ge1$.
Write $P_1=\mathrm{primePairSetFirst}(z)$ and
$P_2(q)=\mathrm{primeSetTwoBounds}(z,z',q)$.
For modulus $8M$, tolerance $\varepsilon/2$ and exceptional fraction
$1/2$, there is $G_1\subset P_1$ and, for each $q\in G_1$,
there is $G_2(q)\subset P_2(q)$ with the prescribed
$\mathrm{GoodQ}$ and $\mathrm{GoodD}$ character boxes and
\[
 |G_1|\ge |P_1|/2>0,\qquad |G_2(q)|\ge |P_2(q)|/2>0 .
\]
The first pool consists of primes $q\equiv3\pmod8$ in $(e^{z/2},e^z]$. The second pool consists of primes $d\equiv7\pmod {16}$ between
$e^{z'/2}$ (strict) and $e^{z'}$, with
$q\prod_{3\le p\le z,\,p\ {\rm prime}}p\mid d+1$. The row and box conversion
on these pools yields $\mathrm{PairGood}(M,\varepsilon,Y,Y^J,q,d)$
for selected good primes, including the primitive inducing characters'
bounds and the real-character zero boxes. See
\nolinkurl{QH.pair_share}, \path{MeanSquareBounds/PairShare.lean:26--32};
\nolinkurl{PPset}, \nolinkurl{oddM}, \path{GL3MeanSquares/SecondEstimateOfIntervalHypothesis.lean:50--53};
\nolinkurl{primeSetTwoBounds}, \path{MeanSquareBounds/PrimeSetTwoBounds.lean:41--53};
\nolinkurl{QH.pair_exists}, \path{MeanSquareBounds/PairExists.lean:9--49};
\nolinkurl{QH.pairEx}, \path{MeanSquareBounds/pairEx.lean:112--169}.

For the same $\psi$, \nolinkurl{fieldChi} supplies $c\ge1,e_0$
so that $e\ge e_0$ and primes $q,d$
with $q\ne3$, $(M,qd)=1$, $q\equiv3\pmod 8$, $d\equiv7\pmod {16}$,
$(-d/q)=1$, and $(-d/p)=1$ for all odd primes $p$ up to a cutoff $b_0\ge4M$
(\nolinkurl{HostHyp}, \path{FrobeniusOrbits/Statement.lean:13--14};
\nolinkurl{fieldChi}, \path{CharacterBoxes/FieldChi.lean:20--26})
give a field frame $H_0$ of order $c2^e$, field rows up to the specified
cutoff, conductor factor $Q_S\mid\operatorname{lcm}(2^{e+2},M)$ with
$Q_S\le4Mc2^e$, kept-character data, the Artin kernel law and a twist
description. Each frame in the order range admits ranks and the host
choices of \decid{D513}--\decid{D515}
(\nolinkurl{admissibleClassesTwistedChar}, \path{CoprimeRanks/AdmChi.lean:11--20}).
Given these choices, the field rows, a good pair and a window at $Y$,
\decid{D511} extends $H_0$ to a good host with coset shares, uniformly
for $s.J,R.o,M\le C_0$ and $Q_S\le C_0R.n$.
The precise divisibility and positivity conditions are those of
\nolinkurl{hostOfField}, \path{CharacterBoxes/HostOfField.lean:8--15}.
\end{formalstatement}

\begin{proof}
Write $a=2+8\varepsilon$, $k=c^{1/a}\ge1$. Choose $e(u)$ with
$u^a\le2^{e(u)}\le2u^a$ for $u\ge1$ by rounding up to a dyadic power.
The window hypothesis supplies $J\ge1$ and a threshold for this
particular $e$. For arbitrarily large $X$ put $Y=X^{1/J}$.
The first cardinality inequality gives a $q\in G_1$; with this $q$
fixed, the second gives $d\in G_2(q)$. The good-pair conclusion
gives the admissible row at $Y$. Applying the window hypothesis to
these primes gives $W\ge Y$, $W^2\le X$ and the asserted Selmer
bound at every $u\in[W,W^2]$. Enlarge the starting $X$ so that
$W\ge\max(1,k^3,k^{400})$. Put $Z=kW$. For $Z\le x\le Z^{3/2}$
we have $x\le W^2$ since $k^3\le W$, and $u=x/k\in[W,W^2]$.
Consequently
\[
 x^a\le c2^{e(u)}\le2x^a,\qquad
 x^{400}=k^{400}u^{400}\le u^{401},\qquad
 \exp(u^\varepsilon)\le\exp(x^\varepsilon).
\]
The good-pair bounds restrict from $X$ to $x$.
The two last inequalities and nonnegativity of each Selmer summand
restrict the bound to rank $\lfloor x^{400}\rfloor$ and budget
$\sqrt{\exp(x^\varepsilon)}$ for every field of type
$(\psi,e(u),q,d)$. Thus $x^{1/J}\le Y\le x$ and
the selected datum exists throughout arbitrarily far $3/2$-windows.
Use all scales admitting such data for this $J$.
See \nolinkurl{exists_pageN}, \nolinkurl{VVle.mono},
\path{CharacterBoxes/ExistsPageN.lean:51--77};
\nolinkurl{winPairs}, \path{CharacterBoxes/WinPairs.lean:5--81};
\nolinkurl{WinDatum}, \path{CharacterBoxes/TwistForm.lean:19--26}.

Fix $C_0\ge\max(1,J,\Lambda.o,8M)$ before selecting a scale. Use the
common thresholds from the field extension, admissible ranks and
$R.t\le x^{400}$. Enlarge $x_0$ so that $x>c2^{e_0}$ and
$Y^{\varepsilon/5}\ge\max(5,4M)$ whenever $x\ge x_0$
(which follows from $Y\ge x^{1/J}$). Now take a member scale
and its witnesses $Y,q,d,e$. The order lower bound gives
$c2^e\ge x^a\ge x>c2^{e_0}$. The pair supplies the required prime
congruences and $(-d/p)=1$ for odd primes $p\le\lfloor
Y^{\varepsilon/5}\rfloor$, with $q\ne3$ by its lower bound (\nolinkurl{pairHyp}, \path{CharacterBoxes/PairHyp.lean:24--31}).
Apply the field input at this cutoff and choose the ranks on its exact
frame. Since $R.n\ge x^a$,
\[
 Q_S\le4Mc2^e\le8Mx^a\le C_0R.n .
\]
All the hypotheses of the field extension now hold. In particular the
two sets selected in that construction have different jobs: its
relation-free second set $T_2$ has
$x^{-\varepsilon}R.n\le |T_2|$ and $4L_w|T_2|\le R.n_0$.
The set $TT$ of translates is then chosen with
$\delta_1R.n/2\le |TT|\le2\delta_1R.n$, the weighted bad-prime bound,
and, for every coset $\gamma+A_0$,
\[
  \big|\,|TT\cap(\gamma+A_0)|-\lambda_\gamma |TT|\,\big|
   \le x^{-1+\varepsilon}|TT|.
\]
Here $\lambda_\gamma$ is $4/(3o+1)$ on $A_0$ and $3/(3o+1)$ elsewhere.
Inserting $TT$ preserves $H_0$ and the selected field. These are the
count estimates used inside \decid{D511}
(\nolinkurl{hof_thick}, \nolinkurl{UpperTheory.T3.hostOfFieldOf},
\path{ShortNorms/hostOfFieldOf.lean:71--95}, \path{ShortNorms/hostOfFieldOf.lean:194--227};
\nolinkurl{topsChoiceH}, \path{ThickFrames/topsChoiceH.lean:30--34};
\nolinkurl{Frame.Shares}, \path{ExplicitFormulaErrorOfHostConditions.lean:135--140}).

Every prime divisor of $Q_S$ divides $2M$ by the displayed lcm divisibility.
The twist description and good pair now give
$|G_\chi(x)|\le x^{1/2+\varepsilon}$ and logarithmic-derivative norm
$\le4\log x$ for every bounded character. Together with the kept and
exceptional-character data and the coset shares these are precisely
$\mathrm{EAtH}$. The Artin kernel identifies $H.E$ as one of the fields
quantified over on the selected window. Since $R.t\le\lfloor
x^{400}\rfloor$, monotonicity gives the second assertion on $H.E$.
These constructions have fixed $C_5=4$ and all other constants before
the choice of $s$. See \nolinkurl{boundedTw}, \path{CharacterBoxes/BoundedTw.lean:45--56};
\nolinkurl{isChiField}, \path{CharacterBoxes/IsChiField.lean:5--22};
\nolinkurl{eventually_rank_le_rpow_of_lt_one_div_fiftytwo},
\path{CharacterBoxes/EventuallyRankLeRpowOfLtOneDivFiftytwo.lean:5--9};
\decid{D510}--\decid{D512}, \nolinkurl{supplyOfAdmOf},
\path{ShortNorms/supplyOfAdmOf.lean:28--57}, \path{ShortNorms/supplyOfAdmOf.lean:68--135}.
\end{proof}

\begin{proposition}[Common family for the request]
Under Green and a witness of the closed-link and definite-side
$\mathrm{HypsAlt}$ in \nodeid{P015}, $\mathrm{TopG}$ holds. Thus if
$\lambda>1$, hosts meet a request $G$ on $\lambda$-windows and membership
in $G.\mathrm{cls}(\lambda)$ implies $\mathrm{HypV}$, then there are real
$e_D,e_L$ and $\varepsilon_1>0$
such that each $0<\varepsilon<\varepsilon_1$ has one family in that class
with tied data on all its sufficiently large window scales.
\end{proposition}
\noindent\nolinkurl{TopG}, \path{CharacterBoxes/ZeroFreeAllOfTopG.lean:12--14};
\nolinkurl{FamilyClassConnected}, \path{CharacterBoxes/DivinNOfHypV.lean:13--15}.

\begin{formalstatement}[Design and closed estimates used below]
On intersecting the arithmetic family class with $G.\mathrm{cls}(\lambda)$,
host existence and the design hypotheses give a design for each prescribed
largeness parameter, including $\mathrm{ChosenE}$, $\mathrm{TuSize}$ and
$\mathrm{DGamSize}$ (\decid{D483}, \nolinkurl{designExists_of_tuSize_aux4n_litThin},
\path{ThirdPlaceSeries/LinkEx.lean:15--37}; \nolinkurl{DesignExists},
\path{ParamsHeightRankConstant.lean:187--190}).
For designs in the underlying class, the following nine properties are
closed in precisely the conditional forms of \decid{D488}:
$N$, division given its local inputs, $\eta$, irreducibility, the
small-level rank equality, the two chosen-design rank statements,
$G$, $S$, and the global assembly implication. Here the symbols denote
the respective hypotheses $hN,hDIV,hEta,hIRR,hZ,hRk,hG,hS,hGL$ of
\nolinkurl{topDesWin}, \path{CharacterBoxes/DesWin.lean:11--22}.
The $N$ property is supplied in the hypotheses. In \hyperref[sub:B05]{B05}
the division input uses the selected rows (\decid{D480}--\decid{D482})
and local admissibility \hyperref[sub:B06]{B06}. The latter takes point,
period and column bounds and served selections, uses design size to get
served families and projects admissibility from the inclusion result
(\decid{D516}--\decid{D518},
\nolinkurl{closed_family_admissible_twist_link_alt},
\path{ThirdPlaceSeries/ClosedFamilyAdmissibleTwistLink.lean:12--31}).
The other closed inputs and the restriction of size estimates to the
intersection are those assembled by \nolinkurl{shape_closed_link_n_of_green},
\path{ThirdPlaceSeries/TopShape.lean:33--55}, especially
\decid{D479}, \decid{D484}--\decid{D486}, \decid{D489}.
\end{formalstatement}

\begin{proof}
Fix the request and the host witnesses. Feed the raw selections of the
given witness into the shape, passing the scalar premises
$\mathrm{True},\mathrm{True}$, the projected definite rows, and the
served-twist condition (\nolinkurl{goal_of_shape_of_hyps},
\path{ThirdPlaceSeries/ShapeOf.lean:24--32}). Use the constants
$(C_{\rm abs},C_0,e_D,e_L)=(2,3,2,0)$ of
\path{ThirdPlaceSeries/TopShape.lean:50--55}.
Intersect the nine closed conditions: take the minimum of their
positive $\varepsilon$ thresholds and, at each admissible $\varepsilon$,
the componentwise supremum of their largeness parameters
(\nolinkurl{Closed.and}, \path{ParamsHeightRankConstant.lean:159--171}).
Take the minimum with the design's positive threshold. At any smaller
positive $\varepsilon$, ask for a design at this common largeness
parameter. Its family $F$ projects both to the original class and to
the request. All nine estimates therefore apply at the same $F,C_H$.

The request supplies $\mathrm{HypV}(F)$, which with $N$, the rank
equality and $\varepsilon>0$ yields $\mathrm{DIVinN}(F)$.
Insert this and $\mathrm{LinkEta}$ into $\mathrm{LinkDIVt}$ to get
$\mathrm{LinkDIV}$. The two rank statements at the chosen design,
together with the size bounds and division, give $\mathrm{LinkMu}$.
The global implication receives the chosen design and the five entries
$(\mathrm{LinkG},\mathrm{LinkS},\mathrm{LinkIRR},\mathrm{LinkDIV},\mathrm{LinkMu})$.
Its conclusion gives $K,c_{HH},e_K,C_L,x_0$ and the tied estimate
for every $x\in F.\mathrm{Win}$ above $x_0$. This is precisely
the connected-family assertion. See \nolinkurl{divinN_of_hypV},
\path{CharacterBoxes/DivinNOfHypV.lean:117--131};
\nolinkurl{LinkDIVt}, \nolinkurl{LinkGLOBAL},
\path{GlobalTheoremInputs.lean:91--95}, \path{GlobalTheoremInputs.lean:158--164};
\nolinkurl{topDesWin}, \path{CharacterBoxes/DesWin.lean:23--31}.

For the application take $\lambda=3/2$, $G=G_\psi$.
On a family in the request class the second component of $G_\psi$ gives
$\mathrm{HypV}$ after $F.x_1+1$; the first gives the character bounds
with constant $4$ at the same indices. Window membership is itself
the first clause of the class. The preceding two propositions thus give
one family carrying the windows, bounds and ties required by the
cyclic-level analytic implication for each $\psi$. Applying that
implication with the level estimate and its $\mathrm{formB}$ input
gives the conditional conclusion of P015. See \nolinkurl{FReq.cls},
\path{CharacterBoxes/Supply.lean:16--23}; \nolinkurl{zeroFreeAll_of_topG},
\path{CharacterBoxes/ZeroFreeAllOfTopG.lean:16--34};
\nolinkurl{zero_free_all_of_green_v_win}, \path{ThirdPlaceSeries/TopPlug.lean:12--15}.
\end{proof}
